%
% 3-risch.tex
%
% (c) 2026 Prof Dr Andreas Müller
%
\section{Der Risch-Algorithmus
\label{buch:intalg:section:risch}}
\kopfrechts{Der Risch-Algorithmus}%
Das Liouville-Prinzip erklärt, wie ein Funktion aussehen muss, wenn sie
eine elementare Stammfunktion haben soll.
Für sich allein genommen reicht es nicht, für eine beliebige elementare
Funktion zu entscheiden, ob sie auch ein elementare Stammfunktion hat.
Dazu sind weitere Überlegungen nötig, die in diesem Kapitel dargestellt
werden sollen.
Das Ziel ist ein Entscheidungsalgorithmus, der ausgehend von einer
elementaren Funktion als Integranden schrittweise eine Stammfunktion
konstruiert.
Falls ein Schritt sich als nicht durchführbar erweist, bedeutet dies,
dass keine elementare Stammfunktion existiert.

Der Risch-Algorithmus entstand aus den Ansätzen, die Richard Henry Risch
in den Sechzigerjahren entwickelte \cite{buch:risch}.
Er beginnt mit einer Verallgemeinerung der Methode von Hermite zur
Behandlung des rationalen Teils.
Die besonders erfolgreiche Methode von Rothstein-Trager zur Bestimmung
der Koeffizienten $c_i$ und der Funktionen $v_i$ kam in den Siebzigerjahren
dazu.
Die folgende Darstellung des Algorithmus stützt sich weitgehend auf 
\cite{buch:algorithms}.

%
% Übersicht über den Algorithmus
%
\subsection{Übersicht über den Algorithmus
\label{buch:intalg:risch:uebersicht}}

\subsubsection{Die Aufgabenstellung}

\subsubsection{Der polynomielle und der rationale Teil}

\subsubsection{Rekursion}

%
% Risch-Algorithmus für logarithmische Erweiterungen
%
\subsection{Risch-Algorithmus für logarithmische Erweiterungen
\label{buch:intalg:risch:logarithmisch}}

\subsubsection{Hermite-Methode}

\subsubsection{Teilbarkeit von $b'(\vartheta)$ und $Db(\vartheta)$}

\subsubsection{Rothstein-Trager für logarithmische Erweiterungen}

\begin{beispiel}
\label{buch:intalg:risch:bsp:1/logx}
Hat die Funktion
\[
f(x)
=
\frac{1}{\log x}
\]
eine elementare Stammfunktion?
Der Integrand liegt in der logarithmischen Erweiterung
$F(\vartheta) = \mathbb{Q}(x,\log x)$ von $F=\mathbb{Q}(x)$
um das logarithmische Element $\vartheta=\log x$.
$f(x)$ hat nur rationalen Teil
\[
f(x)
=
\frac{1}{\vartheta}
\qquad
\Rightarrow
\qquad
\begin{cases}
a(\vartheta) &= 1 \\
b(\vartheta) &= \vartheta.
\end{cases}
\]
Dies ist bereits die Ausgangsform für die Methode von Rothstein-Trager,
die verlangt, dass die Werte von $z$ bestimmt werden müssen, für die
die beiden Polynome
\[
a(\vartheta) - z\cdot Db(\vartheta)
1 - z\,D\vartheta
=
1- z\frac{1}{x}
\qquad\text{und}\qquad
b(\vartheta)=\vartheta
\]
gemeinsame Teiler haben.
Da das linke Polynom Grad 0 hat, kann es nur dann ein Polynom vom Grad
mindestens $1$ als Teiler haben, wenn es verschwindet.
Die Bedinungung für gemeinsame Teiler ist daher
\[
1-z\frac{1}{x} = 0 
\qquad\Rightarrow\qquad
z=x,
\]
insbesondere ist $z$ nicht konstant.
Es folgt, dass $f$ kein elementare Stammfunktion hat.
\end{beispiel}

Im Beispiel wurde bewiesen, dass $1/\log x$ keine elementare Stammfunktion
hat.
Die Funktion hat aber einige Bedeutung in der analytischen Zahlentheorie,
da sie ein Approximation für die Primzahlzählfunktion 
\[
\pi(x)
=
\bigl|
\{ p \mid \text{$p$ ist prim und $p\le x$} \}
\bigr|
\]
ist.
Es bleibt daher nichts anderes übrig, als eine neue spezielle Funktion
dafür zu definieren.

\begin{definition}[Integrallogarithmus]
Die Stammfunktion
\[
\operatorname{li}(x)
=
\int_0^x
\frac{1}{\log t}
\,dt
\]
heisst der \emph{Integrallogarithmus}.
\end{definition}

\begin{beispiel}
Hat die Funktion $1/x\log x \in \mathbb{Q}(x,\log x)$ eine elementare
Stammfunktion?
Mit $\vartheta=\log x$ ist die Stammfunktion des rationalen Teils
\begin{equation}
\frac{1}{x\vartheta}
=
\frac{1/x}{\vartheta}
=
\frac{a(\vartheta)}{b(\vartheta)}
\qquad\Rightarrow\qquad
\begin{cases}
a(\vartheta) &= 1/x\\
b(\vartheta) &= \vartheta
\end{cases}
\label{buch:intalg:risch:eqn:1xlogx}
\end{equation}
nach der Methode von Rothstein-Trager zu bestimmen.
Dazu müssen gemeinsame Teiler von
\[
a(\vartheta) - z\cdot Db(\vartheta)
=
\frac{1}{x} - zD\vartheta
=
\frac{1}{x} - z\frac{1}{x}
\qquad\text{und}\qquad
b(\vartheta)=\vartheta
\]
bestimmt werden.
Da das linke Polynom Grad 0 hat, kann es nur dann nichttriviale
Teiler vom Grad mindestens 1 geben, wenn es verschwindet, wenn also
\[
\frac{1}{x} - z\frac{1}{x}
=
0
\qquad\Rightarrow\qquad
\frac{1-z}{x}
=
0
\qquad\Rightarrow\qquad
c=z=1.
\]
Der gemeinsame Teiler ist $v=\vartheta=\log x$.
Nach der Methode von Rothstein-Trager ist daher
\[
f(x)
=
1\frac{Dv}{v}
=
\frac{1}{x \vartheta}
\qquad\Rightarrow\qquad
\int f
=
\log v
=
\log \vartheta
=
\log\log x,
\]
ein elementare Stammfunktion des Integranden $f$.


Die Wahl $a(\vartheta)=1/x$ in \eqref{buch:intalg:risch:eqn:1xlogx}
wird dadurch festgelegt, dass $b(\vartheta)$ ein normiertes Polynom
sein.
Der Faktor $x$ im Nenner von \eqref{buch:intalg:risch:eqn:1xlogx}
kann daher nicht Teil des Polynoms $b(\vartheta)$ sein, sondern muss
in $a(\vartheta)$ integriert werden.
\end{beispiel}

%
% Der polynomielle Teil für logarithmische Erweiterungen
%
\subsubsection{Der polynomielle Teil für logarithmische Erweiterungen}
Wir betrachten wieder die Situation einer logarithmischen Erweiterung
$F(\vartheta)\supset F$ um ein logarithmisches Element $\vartheta =\log u$
mit Konstantenkörper $K$.
Wir versuchen die Stammfunktion des polynomiellen Teils in $F(\vartheta)$
zu bestimmen, den wir als
\begin{equation}
p(\vartheta)
=
p_0 + p_1\vartheta + \dots + p_{l-1}\vartheta^{l-1} + p_l\vartheta^l
\in
F[\vartheta]
\label{buch:intalg:risch:eqn:logpolyintegrand}
\end{equation}
schreiben.
Für ein Polynom in der unabhängigen Variablen $x$ ist die Stammfunktion
trivialerweise elementar, da 
\[
\int p_kx^k \,dx = \frac{p_k}{k+1}z^{k+1} + C
\]
ist.
Dies funktioniert jedoch nur für konstante Koeffizienten.
In der vorliegenden allgemeinen Situation
\eqref{buch:intalg:risch:eqn:logpolyintegrand}
erwarten wir eine wesentlich kompliziertere Lösung.

Falls $p(\vartheta)$ eine elementare Stammfunktion hat, dann muss 
sie nach dem Liouville-Prinzip die Form
\begin{align}
\int p(\vartheta)
&=
v_0(\vartheta) + \sum_{i=1}^n c_i\log v_i(\vartheta)
\notag
\intertext{mit $v_i(\vartheta) \in F(\vartheta)$ und $c_i\in K$ haben.
Wie früher darf man annehmen, das die $v_i(\vartheta)$ teilerfremde
quadratfreie Polynome sind.
In der abgeleiteten Form ist}
p(\vartheta)
&=
Dv_0(\vartheta)
+
\sum_{i=1}^n c_i \frac{Dv_i(\vartheta)}{v_i(\vartheta)}
\notag
\intertext{Auf der linken Seite dieser Gleichung steht ein Polynom,
auf der rechten Seite kommen die Nenner $v_i(\vartheta)$ vor.
Da sie teilerfremd sind, können sich die Brüche nicht wegheben, in den
Termen der Summe darf also $\vartheta$ gar nicht vorkommen.
Es bleibt daher die Form}
p(\vartheta)
&=
Dq(\vartheta) + \sum_{i=1}^n c_i \frac{Dv_i}{v_i}
\notag
\\
&=
D(q_{l+1}\vartheta^{l+1} + q_l\vartheta^l + \dots + q_2\vartheta^2
+ q_1\vartheta) + Dq_0 + \sum_{i=1}^n c_i\frac{Dv_i}{v_i}
\label{buch:intalg:risch:logarithmisch:eqn:pDq}
\end{align}
mit $v_i\in F$, $c_i\in K$ und $q\in F[\vartheta]$.
Die Aufgabe ist, das Polynom $q(\vartheta)$ zu bestimmen.

Der Koeffizient $q_0$ für sich alleine interessiert nicht wirklich,
da für die gesuchte Stammfunktion nur die Kombination
\[
\bar{q}_0
=
q_0 + \sum_{i=1}^n c_i\log v_i
\]
benötigt wird.

Zur Bestimmung der Koeffizienten $q_k$ und $\bar{q}_0$ vergleichen wir
Koeffizienten in 
\eqref{buch:intalg:risch:logarithmisch:eqn:pDq}.
Es ergeben sich die Gleichungen
\begin{align*}
0   &= Dq_{l+1}
	&&\Rightarrow
	& Dq_{l+1} &= 0
\\[-4pt]
p_l &= Dq_l + q_{l+1}(l+1) \frac{Du}{u}
	&&\Rightarrow
	& Dq_l     &= p_l - (l+1)q_{l+1}\frac{Du}{u}
\\[-9pt]
    &\phantom{n}\vdots 
	&&           &
	&\phantom{n}\vdots
\\[-9pt]
p_k &= Dq_k + q_{k+1}(k+1) \frac{Du}{u}
	&&\Rightarrow
	& Dq_k     &= p_k - (k+1)q_{k+1}\frac{Du}{u}
\\[-9pt]
    &\phantom{n}\vdots        
	&&
	&          &\phantom{n}\vdots
\\[-9pt]
p_1 &= Dq_1 + 2q_2 \frac{Du}{u}
	&&\Rightarrow
	& Dq_1     &= p_1 - 2q_2 \frac{Du}{u}
\\[-0pt]
p_0 &= D\bar{q}_0 + q_1 \frac{Du}{u}
	&&\Rightarrow& D\bar{q}_0 &= p_1 - q_1 \frac{Du}{u}
\end{align*}
Die rechten Seiten zeigen, dass alle Koeffizieten von $q(\vartheta)$
als Lösung eines Integrationsproblems im kleineren Körper $F$ gefunden
werden müssen.
Jedes dieser einzelnen Integrationsproblem muss eine elementare
Lösung in $F(\vartheta)$ haben.
Die Lösung des Integrationsproblems für $d_k$ ist aber immer nur bis auf
eine Integrationskonstante $b_k$ bestimmt.

Für die erste Gleichung folgt, dass $q_{l+1}$ eine Konstante sein muss,
die mit $b_{l+1}$ bezeichnet wird.
Da $q_{l+1}$ auch in der Gleichung für für $q_l$ auftritt, kann
letztere zur Bestimmung von $b_{l+1}$ herangezogen werden.

Die Gleichung für $q_l$ wird mit der Lösung für $q_{l+1}$ zu
\[
Dq_l
=
p_l-(l+1)b_{l+1}\frac{Du}{u}.
\]
Die Lösung muss eine elementare Funktion sein, nach dem Prinzip von Liouville
muss
\[
q_l
=
d_l + c_l\vartheta + b_l
\]
mit $d_l\in F$ und $c_l,b_l\in K$ sein.
Durch Ableiten und Vergleich mit der Gleichung für $Dq_l$ schilessen wir,
dass
\[
p_l - (l+1)b_{l+1}\frac{Du}{u} 
=
Dq_l 
=
Dd_l + c_l\frac{Du}{u}
\qquad\Rightarrow\qquad
\left\{
\;
\renewcommand{\arraycolsep}{2pt}
\begin{array}{rcl}
\displaystyle d_l     &=& \displaystyle \int p_l \\[9pt]
\displaystyle b_{l+1} &=& \displaystyle -\frac{c_{l}}{l+1}.
\end{array}
\right.
\]
sein muss.
Damit ist jetzt auch die Integrationskonstante $b_{l+1}$ bestimmt.

Nehmen wir an, dass die Koeffizienten bis $q_{k+1} = d_{k+1}+b_{k+1}$
bereits bestimmt sind.
Nur die Integrationskonstante $b_{k+1}$ ist noch zu bestimmen.
Die Bestimmungsgleichung
\begin{align*}
Dq_k &= p_k - (k+1) q_{k+1} \frac{Du}{u}
\intertext{wird damit zu}
Dq_k
&=
p_k
-
(k+1) d_{k+1} \frac{Du}{u}
-
(k+1) b_{k+1} \frac{Du}{u}.
\intertext{Der letzte Term kann ist sofort integrierbar, wir können
daher die Terme so sortieren, dass}
p_k
-
(k+1) d_{k+1} \frac{Du}{u}
&=
Dq_k
+
(k+1) b_{k+1} \frac{Du}{u}.
\end{align*}
Die linke Seite muss also eine elementare Stammfunktion haben,
die nach dem Prinzip von Liouville wieder in der Form $d_k + c_k\vartheta$
geschrieben werden kann.
Durch Ableiten und Termvergleich können wir wieder
\[
\left.
\renewcommand{\arraycolsep}{2pt}
\begin{array}{rcl}
d_k
&=&
\displaystyle
\int\biggl(
p_k
-
(k+1) d_{k+1} \frac{Du}{u}
\biggr)
\\[9pt]
b_{k+1}
&=&
\displaystyle
-
\frac{c_k}{k+1}
\end{array}
\;
\right\}
\quad\Rightarrow\quad
q_k = d_k + b_k
\]
schliessen.
Dieser Schritt kann für kleiner werdende $k$ wiederholt werden bis
$q_1=d_1+b_1$ bsetimmt ist, wobei die Integrationskonstante $b_1$ noch
zu finden ist.

\input{chapters/070-intalg/fig/polylog.tex}%

Die Bestimmungsgleichung für $\bar{q}_0$ ist
\[
D\bar{q}_0
=
p_0 - q_1\frac{Du}{u}
=
p_0 - d_1\frac{Du}{u} - c_1\frac{Du}{u}
\quad\Rightarrow\quad
D\bar{q}_0
-
c_1\frac{Du}{u}
=
p_0 - d_1\frac{Du}{u}.
\]
Um $\bar{q}_0$ zu bestimmen, muss eine Stammfunktion gefunden werden.
Dazu muss wieder das Integrationsproblem für den Integranden ganz rechts
gelöst werden.
Im Gegensatz zu den früheren Schritten darf die Stammfunktion
\[
\int
\biggl(
p_0 - d_1\frac{Du}{u}
\biggr)
=
d_0
=
v_0 + \sum_{i=1}^m \bar{c}_i \log v_i, \qquad \bar{c}_i\in K
\]
aber auch andere Logarithmen als $\vartheta$ enthalten, solange
sie Logarithmen von Elemente von $F$ sind.
Einer der Logarithmen darf $\vartheta$ sein, sei dies der erste.
Aus der Konstanten $\bar{c_1}$ lässt sich dann auch auch
\(
b_1 = -\bar{c}_1
\)
bestimmen.
Dadurch ist $\bar{q}_0=d_0$ jetzt bis auf eine Konstante $b_0$ bestimmt,
die nicht weiter bestimmt werden kann.

Der Gang der Berechnung ist in
Abbildung~\ref{buch:intalg:risch:logarithmisch:fig:polylog} 
schematisch dargestellt.
Im Schritt zur Bestimmung von $q_k$ wird zunächst aus dem bekannten
Koeffizienten $p_k$ und der Stammfunktion $d_{k-1}$ aus dem vorangegangenen
Schritt der Integrand gebildet.
Seine Stammfunktion liefert $d_k$ und die Konsante $c_k$.
Die Integrationskonstante $b_{k-1}$ im vorangegangenen Schritt ist
durch $c_k$ bestimmt.
Schliesslich kann $q_k = d_k+b_k$ gefunden werden.
Nur die letzte Integrationskonstante $b_0$ kann nicht bestimmt werden.

\begin{beispiel}
Man finde die Stammfunktion von $f(x) = \log x$.
Der Integrand ist das Polynom
\[
f(x)=\vartheta
=
1\cdot\vartheta + 0
=
p_1\vartheta + p_0
\qquad\Rightarrow\qquad
\begin{cases}
p_1 &= 1\\
p_0 &= 0.
\end{cases}
\]
in der Erweiterung
$F(\vartheta)=\mathbb{Q}(x,\vartheta)$ von $F=\mathbb{Q}(x)$ um das
logarithmische Element $\vartheta=\log x$.
Es muss also ein Polynom
\[
q(\vartheta)
=
q_2\vartheta^2 + q_1\vartheta + \bar{q}_0
\]
vom Grad 2 gefunden werden.
\begin{itemize}
\item Fall $k=2$:
$q_2$ muss eine Konstante $b_2$ sein, deren Wert im
nächsten Schritt bestimmt wird.
\item Fall $k=1$:
Zur Bestimmung von $q_1$ muss zunächst die Stammfunktion
von $p_1$ bestimmt werden, sie ist
\[
\int p_1
=
\int 1
=
x + 0\cdot\vartheta
\quad
\Rightarrow
\quad
\left\{
\quad
\renewcommand{\arraycolsep}{2pt}
\begin{array}{rcl}
b_2 &=& \displaystyle -\frac{c_1}{2} \\
d_1 &=& x
\end{array}
\quad
\right\}
\quad
\Rightarrow
\quad
q_1 = x + b_1
\]
Da darin $\vartheta$ nicht vorkommt, folgt $c_1=0$ und daher $b_2=0$.
Die Integrationskonstante $b_1$ wird im nächsten Schritt bestimmt.
\item Fall $k=0$:
Es muss die Stammfunktion von $p_0-d_1\,Du/u$ aus
\[
q_0
=
\int\biggl(p_0-d_1\frac{Du}u\biggr)
=
\int\biggl(0-x\frac{1}{x}\biggr)
=
-
\int 1
=
-x
+
b_0
=
d_0+b_0
\]
bestimmt werden.
Auch hier kann man schliessen, dass $b_1=0$ ist.
\end{itemize}
Aus den drei Fällen lässt sich jetzt die Stammfunktion
\[
\int p(\vartheta)
=
q(\vartheta)
=
q_0 + q_1\vartheta + q_2\vartheta^2
=
-x+b_0 + z\vartheta
=
x\log x - x + b_0
\]
ablesen.
Tatsächlich ist die Ableitung davon
\[
\frac{d}{dx}\bigl( x\log x - x + b_0\bigr)
=
\log x + x \frac 1x - 1
=
\log x.
\]
Wie erwartet kann die Integrationskonstante $b_0$ nicht bestimmt werden.
\end{beispiel}

\begin{beispiel}
Hat der Integrand $\log\log x$ eine elementare Stammfunktion?
Für diese Aufgabe muss der Körper der rationalen Funktionen zweimal
logarithmisch erweitert werden, zunächst um $\vartheta_1=\log x$ und
dann im zweiten Schritt um $\vartheta_2 = \vartheta =\log u = \log\log x$.
Gesucht ist jetzt das Integral des Polynoms
$p(\vartheta) = 1\cdot\vartheta + 0$
mit $p_1=1$ und $p_0=0$.
\begin{itemize}
\item Fall $k=2$:
$q_2$ muss eine Konstante $b_2$ sein, deren Wert im
nächsten Schritt bestimmt wird.
\item Fall $k=1$:
Die Stammfunktion $p_1$ ist
\[
\int p_1
=
\int 1
=
x + 0\cdot \vartheta
\quad
\Rightarrow
\quad
\left\{
\quad
\renewcommand{\arraycolsep}{2pt}
\begin{array}{rcl}
b_2 &=& \displaystyle -\frac{c_1}{2} \\
d_1 &=& x
\end{array}
\quad
\right\}
\quad
\Rightarrow
\quad
q_1 = x + b_1
\]
Wie früher folgt $c_1=0$ und $b_2=0$.
Die Integrationskonstante $b_1$ wird im nächsten Schritt bestimmt.
\item Fall $k=0$: Es ist die Stammfunktion von $p_0-d_1 \, Du/u$ zu bestimmen:
\[
q_0
=
\int\biggl(p_0-d_1\frac{Du}u\biggr)
=
\int\biggl(0-x \frac{D\log x}{\log x}\biggr)
=
-\int x \frac{1/x}{\log x}
=
-
\int 
\frac{1}{\log x}.
\]
Dieses Integral ist nach
Beispiel~\ref{buch:intalg:risch:bsp:1/logx}
nicht elementar und daher kann auch das Integral von $\log\log x$ nicht
elementar sein.
\qedhere
\end{itemize}
\end{beispiel}

%
% Risch-Algorithmus für exponentielle Erweiterungen
%
\subsection{Risch-Algorithmus für exponentielle Erweiterungen
\label{buch:intalg:risch:exponentiell}}

\subsubsection{Hermite-Methode}

\subsubsection{Rothstein-Trager für exponentielle Erweiterungen}

\subsubsection{Erweiterte Polynome}

\subsubsection{Der polynomielle Teil und die Risch-Differentialgleichung}

%
% Risch-Algorithmus für algebraische Erweiterungen
%
\subsection{Risch-Algorithmus für algebraische Erweiterungen
\label{buch:intalg:risch:algebraisch}}







